\section{Results}
\label{sec:results}
\begin{table*}[t]
\centering
\caption{Localization of real logical defects (AU-sPRO up to 5\% false-positive rate, \%). Class-wide and footprint heads are retrained from the released head on identical inputs; the last two columns are paired differences with 95\% bootstrap intervals. Bold marks the best head in each row.}
\label{tab:native}
\small\setlength{\tabcolsep}{4.6pt}
\begin{tabular}{llrrrrrcc}
\toprule
\multicolumn{2}{l}{Native test set} & Released & \multicolumn{2}{c}{Class-wide} & \multicolumn{2}{c}{Footprint} & \multicolumn{2}{c}{Footprint $-$ class-wide} \\
\cmidrule(lr){4-5}\cmidrule(lr){6-7}\cmidrule(lr){8-9}
 & & & seed 1 & seed 2 & seed 1 & seed 2 & seed 1 & seed 2 \\
\midrule
Pushpins & all logical & 33.9 & 34.9 & 33.4 & \textbf{78.5} & 71.6 & +43.6 [26.6, 56.0] & +38.2 [16.5, 54.7] \\
 & extra pin & 16.4 & 14.0 & 18.7 & \textbf{83.7} & 70.1 & +69.7 [47.4, 79.7] & +51.4 [21.0, 69.1] \\
 & missing pin & 45.2 & 47.7 & 41.7 & \textbf{70.8} & 69.6 & +23.2 [5.7, 41.9] & +27.9 [9.4, 48.5] \\
 & missing separator & 72.3 & 85.6 & 77.0 & 94.0 & \textbf{95.3} & +8.4 [5.6, 11.4] & +18.3 [10.9, 25.1] \\
\midrule
Breakfast box & all logical & 18.3 & 18.2 & 18.5 & 27.2 & \textbf{30.7} & +9.0 [4.2, 16.2] & +12.1 [6.9, 20.2] \\
 & missing cereals & 55.1 & 60.0 & 54.7 & 86.3 & \textbf{89.2} & +26.4 [11.6, 43.3] & +34.5 [18.3, 52.7] \\
\midrule
Juice bottle & all logical & 33.6 & \textbf{39.5} & 38.3 & 39.4 & 38.4 & $-$0.1 [$-$1.4, 1.9] & +0.0 [$-$1.5, 2.7] \\
\midrule
Splicing conn. & all logical & 17.8 & 18.0 & \textbf{23.6} & 18.8 & 18.5 & +0.8 [$-$0.7, 4.8] & $-$5.1 [$-$11.9, 1.0] \\
\midrule
Screw bag & all logical & \textbf{17.6} & 15.1 & 14.4 & 15.4 & 15.6 & +0.3 [$-$5.0, 2.6] & +1.2 [$-$2.1, 2.4] \\
 & extra part & 2.5 & \textbf{4.3} & 0.6 & 0.0 & 0.0 & $-$4.3 [$-$16.8, $-$0.1] & $-$0.6 [$-$9.6, $-$0.0] \\
\bottomrule
\end{tabular}

\end{table*}

\begin{figure*}[t]
\centering
\includegraphics[width=\linewidth]{native_results.pdf}
\caption{Real logical defects. (a) sPRO against false-positive rate on pushpins; dotted lines are the second seed. (b) AU-sPRO by defect type; bars are seed means, dots are seeds with 95\% intervals. (c) Defective compartments found and intact compartments opened on real pushpin defects, at the scene cutoff, threshold and list length fixed on development trays; whiskers are 95\% intervals.}
\label{fig:native}
\end{figure*}

\subsection{Footprint labels double localization of real defects}
\label{sec:native}
\Cref{tab:native} and \cref{fig:native} give our main comparison. On the 91 real logical pushpin anomalies, footprint labels raise AU-sPRO from \ppCW\% for the matched class-wide heads (\ppRel\% for the released head) to \ppFP\%, a paired gain of \ppGain{} points in the two seeds. The gain is largest where \cref{prop:optimal}(i) predicts it. On extra-pin defects, where a class-wide map lights up every pin, AU-sPRO rises by \ppAddGain{} points. Missing pins and missing separators improve as well. A vacated pin or separator turns into tray floor, and a class-wide label for an inserted floor region covers the entire floor, so the class-wide maps flag most of the tray (\cref{fig:examples}). \Cref{fig:native}a shows that the footprint curves dominate over the whole low false-positive range. The result does not depend on the training recipe: across all six matched pushpin pairs from three recipes, including two earlier ones in the supplement, footprint labels reach \sixFP\% and class-wide labels \sixCW\%.

Breakfast boxes confirm the effect in a different scene. Footprint labels raise AU-sPRO from \bbCW\% to \bbFP\% (\bbGain{} points), driven by missing cereals (\bbCerealCW\% to \bbCerealFP\%), a fault confined to one compartment. In juice bottles, whose labels and liquid are single instances, an inserted part's class channel and its footprint nearly coincide, and AU-sPRO is unchanged (\jbFP\% against \jbCW\%; \jbGain{} points). Splicing connectors show no consistent change either (\scGain{} points). Image-level detection is essentially preserved: AUROC is \bbAucFP\% against \bbAucCW\% in breakfast boxes, and \aucFP\% against \aucCW\% in pushpins, where both continuations remain below the released head (\aucRel\%) that was trained jointly for 70{,}000 steps.

\begin{figure}[t]
\centering
\includegraphics[width=\columnwidth]{native_examples.pdf}
\caption{Real pushpin defects (first test image of each type; cyan: annotation). Each map is divided by its head's pixel threshold at 1\% false-positive rate on good test images. For the missing pin and separator, released and class-wide heads flag 64--76\% of non-defect pixels and the footprint head 2--11\%, each covering at least 89\% of the defect.}
\label{fig:examples}
\end{figure}

\subsection{Fewer wasted inspections at fixed operating points}
\label{sec:operating}
\Cref{fig:native}c answers the question an inspector asks, using only cutoffs fixed on development trays. On real pushpin defects, the footprint heads find \opFPR\% of the defective compartments while opening \opFPW{} intact compartments per faulty tray. The released map's threshold finds \opRelThrR\% and opens \opRelThrW. Footprint labels therefore improve the threshold policy in both respects, finding \opVsThrR{} points more defective compartments and opening \opVsThrW{} fewer intact ones. The released map's best list, its top three compartments, opens \opRelListW{} intact compartments but finds only \opRelListR\% of the defective ones. Against this list, footprint thresholds find \opVsListR{} points more defective compartments at a cost of \opVsListW{} extra compartments per faulty tray, and queue fewer compartments on good trays (\opVsListClean). Relative to the matched class-wide heads, footprint thresholds open \opVsCWW{} fewer intact compartments per faulty tray. At the development scene cutoff the footprint heads alarm on \opFPScene\% of the anomalous images and on \opFPGood\% of the good images; the released head alarms on all anomalous images and on \opRelGood\% of good ones.

\begin{table*}[t]
\centering
\caption{Registered replication with SALAD's full 70{,}000-step recipe from random initialization (pushpins). Matched trays: spread (\%), wasted inspections $W$ and target recall $R$ (\%) at the development-fixed threshold, and pooled $W$. Real defects: AU-sPRO up to 5\% FPR and on extra pins (\%), defective compartments found (\%) and intact compartments opened per faulty tray at the development-fixed operating point, and image AUROC (\%). Bold marks the best retrained head.}
\label{tab:replication}
\small\setlength{\tabcolsep}{5pt}
\begin{tabular}{lrrrrrrrrr}
\toprule
Head & \multicolumn{4}{c}{Matched trays} & \multicolumn{5}{c}{Real pushpin defects} \\
\cmidrule(lr){2-5}\cmidrule(lr){6-10}
 & Spread & $C$ & $R$ & Pooled $C$ & AU-sPRO & Extra pin & Found & Opened & AUROC \\
\midrule
Class-wide, seed 1 & 25.4 & 12.02 & 96.0 & 8.95 & 29.7 & 23.8 & 77.6 & 6.12 & 96.6 \\
Class-wide, seed 2 & 97.9 & 11.45 & \textbf{96.7} & 4.92 & 44.1 & 23.9 & 68.8 & 5.36 & \textbf{99.8} \\
Footprint, seed 1 & 1.8 & \textbf{0.24} & \textbf{96.7} & \textbf{0.21} & \textbf{93.4} & \textbf{90.0} & \textbf{79.9} & \textbf{0.11} & 99.6 \\
Footprint, seed 2 & \textbf{1.6} & 0.59 & 94.6 & 0.48 & 81.7 & 84.1 & 76.2 & 0.35 & 96.7 \\
\midrule
Released (reference) & 89.4 & 11.07 & 95.7 & 4.58 & 33.9 & 16.4 & 51.0 & 5.27 & 99.7 \\
\bottomrule
\end{tabular}

\end{table*}

\subsection{Replication with SALAD's full training recipe}
\label{sec:replication}
The heads above continue the released detector. To test the label under SALAD's complete training, we retrained the composition branch, both the reconstruction autoencoder and the discriminative U-Net, from random initialization with the published 70{,}000-step recipe. We registered the endpoint, the guardrails and a reproduction check before scoring. As before, each seed trains both arms in one process from the same initialization and the same edit stream. \Cref{tab:replication} gives the result for both seeds. The class-wide arms detect every fault and open as many intact compartments as the released head (\scrThrWCW{} per faulty tray at the threshold). In seed 2 the class-wide arm reproduces the released spread closely (97.9\%). In seed 1 it spreads less (25.4\%) but ranks the targets worse, failing our registered reproduction check (supplement). The footprint arms remove the spread (\scrSpreadFP\%). At the 95\%-recall threshold it opens \scrThrWFP{} intact compartments per faulty tray against \scrThrWCW{} for its class-wide twin (\scrDThrW), which also beats the released map's best list (\relListW). On real pushpin defects its AU-sPRO is \scrSproFP\% against \scrSproCW\% (\scrGain{} points). At the development-fixed operating point it finds \scrOpR\% of the defective compartments while opening \scrOpW{} intact compartments per faulty tray, against \opRelListR\% and \opRelListW{} for the released map's best list (\scrVsListR{} points found, \scrVsListW{} compartments opened). Trained in full, footprint labels thus outperform the released detector's best policy on real defects in both respects, in both seeds. Image AUROC shows no consistent cost (\scrAucFP\% against \scrAucCW\% for the class-wide twins). The cost that remains is valid-sham alarms at the development scene cutoff: \scrShamFP{} of 138, against \scrSrcFP{} unedited sources. In seed 1 this exceeds the class-wide twin by more than our registered guardrail allows.

\begin{table*}[t]
\centering
\caption{Matched pushpin trays, re-segmented path, 138 photographs. Alarms out of 138; spread in \%; $W$ is intact compartments opened per faulty tray and $R$ target recall (\%) at the development-fixed threshold or list; pooled $W$ mixes relocations and vacancies; region area is the cell analogue of low-FPR sPRO (\%) with faulty-tray negatives only or with source and sham compartments added. Bold marks the best SALAD head; the rule needs task knowledge.}
\label{tab:matched}
\small\setlength{\tabcolsep}{4.2pt}
\begin{tabular}{llrrrrrrrrrr}
\toprule
\multicolumn{2}{l}{Detector, insertion label} & \multicolumn{2}{c}{Alarms} & Spread & \multicolumn{2}{c}{Threshold} & \multicolumn{2}{c}{List} & Pooled & \multicolumn{2}{c}{Region area} \\
\cmidrule(lr){3-4}\cmidrule(lr){6-7}\cmidrule(lr){8-9}\cmidrule(lr){11-12}
 & & fault & sham & (\%) & $W$ & $R$ & $W$ & $R$ & $W$ & fault & all \\
\midrule
SALAD & released & 138 & 7 & 89.4 & 11.07 & 95.7 & 1.06 & 97.1 & 4.58 & 74.0 & 81.6 \\
 & class-wide, seed 1 & 138 & 19 & 33.8 & 7.72 & 93.5 & \textbf{1.04} & 97.8 & 2.85 & 74.9 & 83.0 \\
 & class-wide, seed 2 & 138 & 15 & 66.5 & 10.87 & 95.3 & 4.04 & 97.8 & 5.34 & 56.8 & 62.8 \\
 & footprint, seed 1 & 138 & 25 & 1.6 & \textbf{4.57} & 96.0 & 2.06 & 97.1 & \textbf{1.09} & 76.2 & \textbf{87.3} \\
 & footprint, seed 2 & 138 & \textbf{6} & \textbf{1.4} & 6.15 & 96.4 & 3.07 & 96.7 & 3.36 & 69.9 & 81.8 \\
\midrule
EfficientAD & -- & 138 & 128 & -- & 0.21 & 97.1 & 0.04 & 97.8 & 0.22 & 92.9 & 69.4 \\
PatchCore & -- & 133 & 123 & -- & 0.82 & 92.8 & 1.96 & 94.9 & 12.50 & 84.3 & 66.4 \\
\midrule
Explicit rule & (reference) & 138 & 0 & -- & 0.00 & 100.0 & -- & -- & 0.00 & -- & -- \\
\bottomrule
\end{tabular}

\end{table*}

\subsection{The mechanism on matched trays}
\label{sec:diagnosis}
The matched trays show why. The released head alarms on all \relFaults{} faulty trays, scores every fault above its valid sham, and ranks the two targets at \relAUCx\% AUROC, with both in its top three on \relTopThree\% of trays (\cref{tab:matched}). Its threshold nevertheless opens \relThrW{} of the 13 intact compartments per faulty tray. The cause is spread, and the controls locate it (\cref{fig:spread}, left). After a relocation, \relAnyFaultX\% of the untouched compartments change by at least 0.5. After the vacancy-only control, which empties the same compartment, the figure is \relAnyVacX\%. After the valid sham, which moves the same pin the same distance, it is \relAnyShamX\%. The response follows the inserted pin. An insertion probe on part maps (\cref{fig:spread}, right) shows that it follows the inserted class. One extra pin raises \ppAny\% of the untouched pins and an extra screw, nut or washer raises \sbAny\% of the untouched parts, while other part classes stay still.

\begin{figure}[t]
\centering
\includegraphics[width=\columnwidth]{spread.pdf}
\caption{Spread. Left: untouched compartments whose map changes by more than a cutoff for each matched edit, with 95\% bands. Right: untouched instances of the inserted class that change by at least 0.5 after one extra instance is pasted into a part map.}
\label{fig:spread}
\end{figure}

Footprint labels remove the spread. With identical inputs, untouched compartments affected by a relocation fall from \cwSpread\% under class-wide labels to \fpSpread\% under footprint labels (paired reductions of \dSpread{} points). In the insertion probe the affected pins fall from \insPPCW\% to \insPPFP\%, and the affected screw-bag parts from \insSBCW\% to \insSBFP\%. Threshold waste on matched relocations falls by \dThrW{} compartments relative to the class-wide pairs. These short continuations leave the targets less separated (both in the top three on \fpTopThree\% of trays), so on matched relocations alone the released map's top-three list (\relListW) stays cheaper than their threshold; full training removes this gap (\cref{sec:replication}).

\subsection{When the faulty location is ambiguous}
\label{sec:ambiguous}
Screw bags are SALAD's motivating case. The rule counts interchangeable screws, nuts and washers, so no single instance is ``the'' extra one. MVTec LOCO annotates them accordingly: for an extra nut the annotation covers every nut in the bag and saturates once one is found, whereas for an extra pushpin it covers only the pins in the overfilled compartment. \Cref{prop:optimal}(ii) predicts that footprint labels cannot help here, and they do not. AU-sPRO is \sbFP\% against \sbCW\% over all screw-bag anomalies (\sbGain{} points). On extra-part defects every head localizes poorly, and the footprint heads lowest (\sbExtraFP\% against \sbExtraCW\%). The mechanism is unchanged, since footprint labels still remove the same-class spread in screw bags. What changes is the task's own definition of where the fault is. The label should follow that definition.

\subsection{What matched shams reveal}
\label{sec:system}
EfficientAD and PatchCore read photographs rather than part maps. Judged on faulty trays alone, they localize the targets better than SALAD's composition head: their fault-only region areas are \eadRegFault\% and \pcRegFault\% against \relRegFault\%, and EfficientAD's threshold opens only \eadThrW{} intact compartments (\cref{tab:matched}). The valid shams overturn this picture. EfficientAD alarms on \eadSham{} and PatchCore on \pcSham{} of 138 valid shams, against \eadSrc{} and \pcSrc{} unedited sources, and both queue the sham's edited compartment. Counting source and sham compartments as negatives reverses the ranking (\eadRegAll\% and \pcRegAll\% against \relRegAll\%). These maps find where pixels were composited, which on faulty trays coincides with the targets. A fault-only evaluation, including sPRO on edited images, cannot separate the two.

\begin{table}[t]
\centering
\caption{SALAD's full system on matched pushpin trays (re-segmented path; alarms out of 138). Native: logical AUROC (\%). The validation cutoff uses SALAD's validation-based normalization with our 10\% split-conformal cutoff; the normal-test cutoff is refit on one half of the unedited test trays and applied to the other.}
\label{tab:system}
\small\setlength{\tabcolsep}{1.9pt}
\begin{tabular}{lrrrrrrr}
\toprule
Branch & Native & \multicolumn{3}{c}{Validation cutoff} & \multicolumn{3}{c}{Normal-test cutoff} \\
\cmidrule(lr){3-5}\cmidrule(lr){6-8}
 & AUROC & src. & sham & fault & src. & sham & fault \\
\midrule
Fused & 100.0 & 75 & 130 & 138 & 13 & 64 & 138 \\
Local app. & 98.5 & 9 & 128 & 138 & 11 & 124 & 138 \\
Composition & 99.7 & 5 & 7 & 138 & 11 & 12 & 138 \\
Gaussian & 84.4 & 84 & 102 & 106 & 8 & 18 & 20 \\
\bottomrule
\end{tabular}

\end{table}

SALAD's full system adds the EfficientAD appearance branch and a global Gaussian branch over pooled student features \cite{rippel2021} and sums their normalized scores with the composition score \cite{fucka2025}. It reaches 100\% logical AUROC on native pushpins. With SALAD's validation-based normalization and a 10\% cutoff on validation scores, however, the fused score alarms on \fullSrcVal{} of 138 unedited test trays (\cref{tab:system}). The Gaussian branch is responsible: its median normalized score is near zero on validation but 2.7 on unedited test trays. Refitting the cutoff on held-out normal test trays restores \fullSrcX{} source alarms. The valid shams then expose a second effect: the fused score still alarms on \fullShamX{} valid shams, because the appearance branch fires on the composite. The composition branch, which reads only the part map, alarms on 12 shams and 11 sources.
